% Preface supplement: Twisted Holography and Quantum Gravity programme.
% Inserted near the end of the preface chapter via \input.

\bigskip
\begin{center}
\rule{0.6\textwidth}{0.4pt}
\end{center}
\bigskip

\section*{IV.\quad Twisted holography and quantum gravity}

\noindent
The twelve sections above construct the universal Maurer--Cartan
element~$\Theta_\cA$ and its projections for every standard
chiral algebra.  This final section asks what these constructions
say about gravity.

A three-dimensional holomorphic-topological (HT) field theory~$T$
on $\R_+\times\C$ is a functor from the bordism category of
surfaces with complex structure and intervals with orientation to
cochain complexes: it assigns a chiral algebra~$\cA$ to the
boundary $\{0\}\times\C$, a Koszul dual~$\cA^!$ governing line
operators, and a line-operator category~$\cC$ to the interval
$\R_+\times\{0\}$.  The holographic principle states that bulk
correlation functions are determined by boundary data, that the
interior is recoverable from the edge.  The question is: what
organises this recovery?

The answer developed in this monograph is that the full
holographic content of~$T$ is packaged into a single algebraic
object (the \emph{holographic modular Koszul datum}) and that
every physical observable of the bulk theory is a projection
of the boundary MC element~$\Theta_\cA$.  The gravitational programme
does not add new axioms to the algebraic engine; it extracts new
consequences from the proved core.  The genus expansion of
$\Theta_\cA$ is the algebraic form of the perturbative expansion
of 3d quantum gravity (after identification of the genus
parameter with the gravitational coupling), the shadow obstruction
tower is the complexity hierarchy of gravitational interactions,
and the Koszul involution $\cA\leftrightarrow\cA^!$ implements
the algebraic counterpart of $S$-duality.

% ====================================================================
\subsection*{The holographic modular Koszul datum}
% ====================================================================

Let $T$ be an HT field theory on $\R_+\times\C$ with boundary
chiral algebra~$\cA$.  The \emph{holographic modular Koszul
datum} of~$T$ is the sextuple
\begin{equation}\label{eq:holo-datum-preface-supp}
\mathcal H(T)
\;=\;
\bigl(
\cA,\;\cA^!,\;\cC,\;r(z),\;\Theta_\cA,\;\nabla^{\mathrm{hol}}
\bigr),
\end{equation}
where:
\begin{itemize}
\item $\cA$ is the boundary chiral algebra on~$\C$.  It is the
  primary algebraic input: every other entry is derived
  from~$\cA$.
\item $\cA^!$ is the chiral Koszul dual, obtained by Verdier
  duality on $\Ran(X)$:
  $\mathbb D_{\Ran}\,\barB_X(\cA)\simeq\cA^!_\infty$.  It is
  the line-side Koszul-dual controller: line operators are
  $\cA^!$-modules, and the bulk algebra is the derived center
  $Z_{\mathrm{der}}(\cA^!)$, one categorical level higher.
\item $\cC$ is the line-operator category.  For a Koszul
  algebra, $\cC\simeq\cA^!\text{-}\mathrm{mod}^{\mathrm{dg}}$:
  the dg derived category of $\cA^!$-modules.  In the
  non-Koszul case, $\cC$ is the category of comodules over
  $\barB(\cA)$, which is the correct derived replacement.
\item $r(z)$ is the \emph{collision $r$-matrix}: the binary
  genus-$0$ component of~$\Theta_\cA$, extracted as the
  collision residue of the two-point shadow at separated points.
  It satisfies the classical Yang--Baxter equation (CYBE) by
  the Arnold relation on $\overline{\FM}_3(\C)$.
\item $\Theta_\cA\in\MC(\gAmod)$ is the universal MC element
  from~\eqref{eq:preface-theta}, proved by the bar-intrinsic
  construction $\Theta_\cA:=D_\cA-d_0$.
\item $\nabla^{\mathrm{hol}}_{g,n}
  =d-\mathrm{Sh}_{g,n}(\Theta_\cA)$ is the \emph{modular shadow
  connection} on the $\overline{\cM}_{g,n}$-family: a flat
  connection whose flatness is the MC equation projected to
  genus~$g$ and arity~$n$.
\end{itemize}

The datum~$\mathcal H(T)$ is \emph{not} a new definition added to
the theory: it is a repackaging of the MC element~$\Theta_\cA$ and
its structural projections into the language of 3d HT field theory.
The five main theorems (A--D+H) are projections of the
scalar-level entry~$\kappa(\cA)$.  The shadow obstruction tower
$\Theta_\cA^{\le r}$ gives the successive nonlinear refinements.
The full datum $\mathcal H(T)$ is the totality.

\subsubsection*{Recovery theorems}

Four recovery theorems establish that the entries of
$\mathcal H(T)$ are not independent but are determined by
the MC element:

\smallskip\noindent
\emph{(R1) Collision residue.}\enspace
$r(z)=\Res^{\mathrm{coll}}_{0,2}(\Theta_\cA)$: the $r$-matrix
is the residue of the genus-$0$, arity-$2$ component of
$\Theta_\cA$ along the collision divisor
$D_{\{1,2\}}\subset\overline{\FM}_2(\C)$.  The codimension-one
residue of the graph-sum formula~\eqref{eq:preface-theta}
extracts the OPE kernel at separated points.

\smallskip\noindent
\emph{(R2) Arnold $\Rightarrow$ CYBE.}\enspace
The Arnold relation
$\eta_{12}\wedge\eta_{23}+\eta_{23}\wedge\eta_{31}
+\eta_{31}\wedge\eta_{12}=0$
on $\overline{\FM}_3(\C)$ implies the classical Yang--Baxter
equation for $r(z)$:
\[
[r_{12}(z_1-z_2),\,r_{13}(z_1-z_3)]
+[r_{12}(z_1-z_2),\,r_{23}(z_2-z_3)]
+[r_{13}(z_1-z_3),\,r_{23}(z_2-z_3)]
\;=\;0.
\]
The three terms of the CYBE correspond to the three terms of
the Arnold relation evaluated on the three codimension-two strata
of $\overline{\FM}_3(\C)$ where all three points collide
simultaneously.  The proof is a direct specialisation of the
$d_{\barB}^2=0$ argument at arity~$3$.

\smallskip\noindent
\emph{(R3) Affine shadow/KZ comparison.}\enspace
On the affine Kac--Moody comparison surface, the modular shadow
connection at genus~$0$ and arity~$n$ is the
Knizhnik--Zamolodchikov connection:
\[
\nabla^{\mathrm{hol}}_{0,n}
\;=\;
d\;-\;\sum_{1\le i<j\le n}r_{ij}(z_i-z_j)\,d\log(z_i-z_j).
\]
Its flatness is the genus-$0$ projection of the MC equation.
The KZ equations, traditionally derived from current-algebra
Ward identities, are here on this affine comparison surface a
structural consequence of $D^2=0$
on $\overline{\FM}_n(\C)$.

\smallskip\noindent
\emph{(R4) Quartic as $L_\infty$-obstruction.}\enspace
The quartic resonance class
$\mathfrak Q(\cA)=[\Theta_\cA^{(4)}]
\in H^2(F^4\gAmod/F^5\gAmod,d_2)$
is the first $L_\infty$-obstruction class of the shadow obstruction tower.
In the holographic interpretation, it is the first bulk
interaction vertex beyond the Gaussian (free) approximation.
Its vanishing ($\mathfrak Q=0$) characterises the algebras
whose bulk theory is one-loop exact at genus~$0$: precisely the
contact/quartic archetype ($\beta\gamma$) and all algebras of
shadow depth~$\le 3$.

These four recovery theorems are not conjectural: they are proved
consequences of the graph-sum formula and the codimension-two
cancellation on~$\overline{\cM}_{g,n}$.

\subsubsection*{The holographic dictionary}

The datum $\mathcal H(T)$ provides a dictionary between algebraic
invariants of~$\cA$ and physical observables of the bulk:
\begin{center}
\small
\renewcommand{\arraystretch}{1.15}
\begin{tabular}{lll}
\textbf{Algebraic entry} & \textbf{Bulk observable}
  & \textbf{Source} \\
\hline
$\kappa(\cA)$ & gravitational central charge
  & genus-$1$ vacuum energy \\[2pt]
$r(z)$ & scattering $S$-matrix at tree level
  & collision residue of $\Theta$ \\[2pt]
$\cA^!$ & line-operator controller
  & Verdier dual \\[2pt]
$\cC\simeq\cA^!\text{-mod}$ & line-operator spectrum
  & Koszul dual modules \\[2pt]
$\nabla^{\mathrm{hol}}_{g,n}$ & gravitational Ward identities
  & shadow connection \\[2pt]
$\mathrm{Sh}_r(\cA)$ & $r$-point bulk vertex
  & arity-$r$ shadow \\[2pt]
$F_g(\cA)$ & genus-$g$ vacuum amplitude
  & scalar trace of $\Theta^{(g)}$ \\[2pt]
$\Delta_\cA(x)$ & spectral branch discriminant
  & $\mathrm{sdet}(1-xT_{\mathrm{br}})$
\end{tabular}
\end{center}

\subsubsection*{Examples: the standard holographic data}

\noindent
\emph{Abelian Chern--Simons} (Heisenberg):
\[
\mathcal H(\cH_k)
\;=\;
\Bigl(
\cH_k,\;\mathrm{Sym}^{\mathrm{ch}}(V^*),\;\Vect,\;\frac{k}{z},\;
k\cdot\eta\otimes\Lambda,\;d
\Bigr).
\]
Every entry is controlled by the single scalar~$k$.  The
$r$-matrix $r(z)=k/z$ is abelian: the CYBE is trivially
satisfied since $[r_{12},r_{13}]=0$.  The shadow connection is
flat with constant curvature.  The line-operator category is
$\Vect$, trivial because the Koszul dual of an abelian algebra
is an abelian coalgebra with no nontrivial modules.  The
holographic datum is maximally simple: one scalar determines the
full gravitational content.

\medskip\noindent
\emph{Non-abelian Chern--Simons}
($\widehat{\mathfrak{sl}}_N$ at level~$k$):
\[
\mathcal H(\widehat{\mathfrak{sl}}_N{}_k)
\;=\;
\Bigl(
\widehat{\mathfrak{sl}}_N{}_k,\;
\widehat{\mathfrak{sl}}_N{}_{-k-2N},\;
\mathcal O_q^{\mathrm{sh}},\;
\frac{\Omega_N}{z},\;
\Theta_k,\;
\nabla^{\mathrm{KZ}}_k
\Bigr),
\]
with $\Omega_N=\sum_a J^a\otimes J_a$ the Casimir,
$q=e^{i\pi/(k+N)}$, $\kappa_k=(N^2-1)(k+N)/(2N)$.  The collision $r$-matrix $r(z)=\Omega_N/z$ satisfies
the CYBE because $\Omega$ is ad-invariant and the triple Casimir
identity holds.  The line-operator category is the
Kazhdan--Lusztig category $\mathcal O_q$ at the corresponding
quantum group parameter; this is the module category for the
dg-shifted Yangian $\Ydg(\mathfrak{sl}_N)$.  The shadow
connection at genus~$0$ and arity~$n$ identifies with the KZ
connection on the affine Kac--Moody comparison surface:
$\nabla=d-\frac{1}{k+N}\sum_{i<j}\Omega_{ij}\,d\log(z_i-z_j)$.

The datum has shadow depth~$3$: the cubic shadow
$\mathfrak C(\widehat{\mathfrak g}_k)$ is nonzero (it detects
the Lie bracket through the triple OPE), but the quartic
resonance class vanishes by cubic gauge triviality
(Theorem~\ref{thm:cubic-gauge-triviality}), since
$H^1(F^3\gAmod/F^4\gAmod,d_2)=0$ for affine Kac--Moody
algebras.  The bulk theory has tree-level interactions but no
irreducible four-point contact vertex.

\medskip\noindent
\emph{3d gravity}
($\mathrm{Vir}_c$):
\begin{align*}
\mathcal H(\mathrm{Vir}_c)
&\;=\;
\Bigl(
\mathrm{Vir}_c,\;
\mathrm{Vir}_{26-c},\;
\cC_c,\;
r_c(z),\;
\Theta_c,\;
\nabla^{\mathrm{hol}}_c
\Bigr),\\[4pt]
r_c(z)
&\;=\;
\frac{c/2}{z^3}+\frac{2T}{z}\,,
\qquad
\kappa_c\;=\;\frac{c}{2}\,.
\end{align*}
The $r$-matrix has a third-order pole (from the fourth-order pole
$T(z)T(w)\sim c/2(z-w)^{-4}+\cdots$ in the OPE; the
$d\log$ extraction absorbs one pole order): it is a
\emph{quasi-trigonometric} solution of the CYBE, not a standard
rational one.  The extra pole orders reflect the conformal-weight
mismatch: $T$ has weight~$2$, so the collision kernel acquires
derivatives.  The shadow obstruction tower is infinite:
$\mathfrak Q^{\mathrm{contact}}_{\mathrm{Vir}}
=10/[c(5c+22)]$, the quintic obstruction is nonzero,
and the holographic datum encodes an infinite sequence of
irreducible bulk vertices.

\medskip\noindent
\emph{Higher-spin gravity}
($\mathcal W_N$ at central charge~$c$):
\[
\mathcal H(\mathcal W_N{}_c)
\;=\;
\Bigl(
\mathcal W_N{}_c,\;
\mathcal W_N{}_{c'_N},\;
\cC_c^{W_N},\;
r_c^{W_N}(z),\;
\Theta_c,\;
\nabla^{\mathrm{hol}}_c
\Bigr),
\]
with $c'_N$ determined by the complementarity
$\kappa_c+\kappa_{c'}=K_N\cdot(H_N-1)$ (a polynomial equation in $c,c'$
depending on $N$; vanishes only for KM, not $\mathcal W_N$).  The $r$-matrix is an
$(N-1)\times(N-1)$ matrix-valued kernel.  The line-operator
category is the module category for the dg-shifted
$\mathcal W_N$-Yangian.  As $N\to\infty$, the datum converges to
the $\mathcal W_\infty$-holographic datum, which governs the
higher-spin limit of 3d gravity.  The completion entropy
$h_K(\mathcal W_N)$ is the kinematic measure of bulk complexity.

\medskip\noindent
\emph{$\beta\gamma$-system} (class~C, shadow depth~$4$):
\begin{align*}
\mathcal H(\beta\gamma)
&\;=\;
\Bigl(
\beta\gamma,\;
(\beta\gamma)^!,\;
\cC_{\beta\gamma},\;
r_{\beta\gamma}(z),\;
\Theta_{\beta\gamma},\;
\nabla^{\mathrm{hol}}_{\beta\gamma}
\Bigr),\\[4pt]
r_{\beta\gamma}(z)
&\;=\;
\frac{\beta\otimes\gamma+\gamma\otimes\beta}{z}\,,
\qquad
\kappa(\beta\gamma)\;=\;c_{\beta\gamma}/2\;=\;6\lambda^2-6\lambda+1.
\end{align*}
The $r$-matrix is symmetric in the $\beta\gamma$-exchange and
has a single simple pole.  The quartic contact invariant
$\mathfrak Q^{\mathrm{contact}}_{\beta\gamma}=0$ vanishes by
the rank-one abelian rigidity
(Corollary~\ref{cor:nms-betagamma-mu-vanishing}), but the full
quartic shadow $\mathrm{Sh}_4\neq 0$ from the log-FM correction.
The Koszul dual $(\beta\gamma)^! = bc$ is a curved chiral algebra
with curvature $m_0=\omega$; the opposite sign $\kappa(bc) = -\kappa(\beta\gamma)$
reflects the duality $c_{bc} = -c_{\beta\gamma}$.
The holographic datum encodes the simplest example with a
nontrivial quartic interaction.

\medskip\noindent
\emph{Lattice vertex algebras} ($V_\Lambda$, class~L for
$\mathrm{rank}\,\Lambda\ge 2$, class~G for rank~$1$):
\[
\mathcal H(V_\Lambda)
\;=\;
\Bigl(
V_\Lambda,\;V_{\Lambda^*}^{\mathrm{tw}},\;
\cC_\Lambda,\;\frac{\Omega_\Lambda}{z},\;
\Theta_\Lambda,\;\nabla^{\mathrm{hol}}_\Lambda
\Bigr),
\]
with $\Omega_\Lambda=\sum_{i,j}\Lambda_{ij}\,\alpha^i
\otimes\alpha^j$ the lattice Casimir (the Gram matrix of the
lattice acting on the Heisenberg currents),
$\kappa(V_\Lambda)=\mathrm{rank}(\Lambda)$ independent of the
cocycle twist (Theorem~\ref{thm:lattice:curvature-braiding-orthogonal}).
The Koszul dual $V_{\Lambda^*}^{\mathrm{tw}}$ is the twisted
lattice algebra on the dual lattice $\Lambda^*=\Hom(\Lambda,\Z)$.
The curvature-braiding orthogonality ensures that the genus-$1$
invariant does not detect the lattice cocycle: the curvature
is a topological invariant of the lattice (its rank), not a
combinatorial one.

\subsubsection*{Functoriality of the holographic datum}

The assignment $\cA\mapsto\mathcal H(\cA)$ is functorial with
respect to the natural operations on chiral algebras:

\smallskip\noindent
\emph{Tensor product.}\enspace
For $\cA\otimes\cA'$ with vanishing mixed OPE:
\[
\mathcal H(\cA\otimes\cA')
\;=\;
\mathcal H(\cA)\times\mathcal H(\cA').
\]
The holographic datum of a product is the product of the data:
each entry decomposes.  In particular, $\kappa$ is additive,
$r(z)=r_\cA(z)\oplus r_{\cA'}(z)$, and the MC element
factorises.  The line-operator category is the tensor product
$\cC\boxtimes\cC'$.

\smallskip\noindent
\emph{Extension of scalars.}\enspace
For a deformation $\cA_t$ of~$\cA=\cA_0$ over a formal parameter
$t$, the holographic datum deforms smoothly:
$\mathcal H(\cA_t)$ is a flat family over $\mathrm{Spec}\,
\Bbbk[\![t]\!]$, with the shadow connection providing the
flat-family structure on $\overline{\cM}_{g,n}$.

\smallskip\noindent
\emph{DS reduction.}\enspace
As detailed below.

\subsubsection*{The holographic morphism under DS reduction}

The Drinfeld--Sokolov reduction $\mathrm{DS}\colon
\widehat{\mathfrak g}_k\to\mathcal W(\mathfrak g,f,k)$
extends to a morphism of holographic data
\[
\mathrm{DS}\colon
\mathcal H(\widehat{\mathfrak g}_k)
\;\longrightarrow\;
\mathcal H(\mathcal W(\mathfrak g,f,k)).
\]
The morphism acts componentwise: the boundary algebra maps by DS,
the Koszul dual maps by the dual DS at level $-k-2h^\vee$, the
$r$-matrix specialises from the Casimir to the $\mathcal W$-matrix
kernel, and the MC element maps by the BRST functor.  The shadow
depth generically increases under DS: for $\mathfrak{sl}_3$
at principal nilpotent, the depth jumps from $3$ (affine) to
$\infty$ (Virasoro/$\mathcal W_3$).  DS reduction creates the
planted-forest shell at genus~$2$ from the tree-level Lie data.
The central charge decomposition $c_{\mathrm{aff}}=c_{\mathcal W}
+c_{\mathrm{ghost}}$ with $c_{\mathrm{ghost}}
=2|\Delta^+(\mathfrak g)|$ is the scalar shadow of the
holographic morphism.

\subsubsection*{The modular shadow connection}

Define the \emph{modular shadow connection} on the universal
curve~$\pi\colon\cC_g\to\overline{\cM}_g$ by
\[
\nabla^{\mathrm{hol}}_{g,n}
\;=\;
d\;-\;\sum_{r\ge 2}\mathrm{Sh}_{g,r}(\Theta_\cA),
\]
where $\mathrm{Sh}_{g,r}$ is the genus-$g$, arity-$r$ shadow
extraction: it evaluates the graph sum at fixed genus, sums over
stable graphs with $r$ external legs, and contracts all internal
data.  The flatness of this connection,
\[
\bigl(\nabla^{\mathrm{hol}}_{g,n}\bigr)^2\;=\;0,
\]
is the MC equation for $\Theta_\cA$ restricted to the
$(g,n)$-stratum.  At genus~$0$ and arity~$n$ it reduces to the
KZ equation.  At genus~$1$ and arity~$0$ it gives
$\kappa(\cA)\cdot\omega_1=0$ on the total space, the curvature
equation.  At genus~$1$ and arity~$1$ it is the heat equation for
conformal blocks with one insertion.

The connection carries a spectral decomposition under the
complementarity splitting.  On the complement
$\cA\oplus\cA^!$, the connection decomposes into the direct-sum
blocks
\begin{align*}
\nabla^{\mathrm{hol}}_\cA
&\;=\;
d-\kappa(\cA)\omega_g
-\mathfrak Q(\cA)\delta_4-\cdots\,,\\[2pt]
\nabla^{\mathrm{hol}}_{\cA^!}
&\;=\;
d-\kappa(\cA^!)\omega_g
-\mathfrak Q(\cA^!)\delta_4-\cdots\,,
\end{align*}
with $\kappa(\cA)+\kappa(\cA^!)=0$ at the scalar level
(for Kac--Moody and free-field algebras; for Virasoro,
$\kappa+\kappa^!=13$).  The
off-diagonal mixing terms (from the complementarity pairing~$C$)
are controlled by the genus-$g$ clutching law.  The full
connection on the complementarity bundle encodes the Ward
identities of the gravitational theory.


% ====================================================================
\subsection*{Perturbative finiteness}
% ====================================================================

\subsubsection*{The HS-sewing criterion}

A chiral algebra~$\cA$ is \emph{perturbatively finite} if its
genus-$g$ partition function $Z_g(\cA)$ converges absolutely for
all~$g$.  The genus expansion of the MC element is a formal power
series in~$\hbar$; its convergence is not guaranteed \emph{a
priori}.  The Hilbert--Schmidt sewing criterion provides the
analytic control.

Let $\cA=\bigoplus_{h\ge 0}\cA_h$ be the conformal-weight
decomposition.  The \emph{sewing kernel} at genus~$g$ is the
bilinear form
\[
K_g(a,b)
\;=\;
\sum_{h,h'\ge 0}
q^{h+h'}\;\langle a_h,\,m_g\,b_{h'}\rangle,
\]
where $m_g$ is the sewing map at genus~$g$ (the contraction over
$g$ pairs of periods on the Schottky uniformisation) and
$q=e^{2\pi i\tau}$ with $\tau$ the modular parameter.  The
kernel $K_g$ is trace-class (and hence $Z_g$ converges) when
$K_g$ satisfies the \emph{Hilbert--Schmidt condition}:
\begin{equation}\label{eq:hs-sewing-preface-supp}
\sum_{h,h'\ge 0}
q^{h+h'}\,\|m_g\|_{h,h'}^2
\;<\;\infty,
\end{equation}
where $\|m_g\|_{h,h'}^2=\sum_{a\in\cA_h,b\in\cA_{h'}}
|\langle a,m_g\,b\rangle|^2$ is the Hilbert--Schmidt norm of
the $(h,h')$-block.

The general HS-sewing theorem (Theorem~\ref{thm:general-hs-sewing}):
\emph{if} $\cA$ has polynomial OPE growth and subexponential
sector growth, \emph{then}~\eqref{eq:hs-sewing-preface-supp}
holds for all~$g$ and all~$q$ in the interior of the Schottky
domain.  Polynomial OPE growth means: the structure constants
$C^c_{ab}$ in the OPE $a(z)b(w)\sim\sum_c C^c_{ab}\,c(w)
/(z-w)^{h_a+h_b-h_c}$ satisfy
$|C^c_{ab}|\le P(h_a,h_b,h_c)$ for a polynomial~$P$.
Subexponential sector growth means: the dimension
$\dim\cA_h=O(e^{C\sqrt{h}})$ for some constant~$C$.  Every
algebra in the standard landscape (Heisenberg, affine
Kac--Moody, $\beta\gamma$, Virasoro, $\mathcal W_N$, lattice
algebras) satisfies both conditions.  The entire standard
landscape is perturbatively finite.

\subsubsection*{The Heisenberg sewing theorem}

For the Heisenberg algebra $\cH_k$, the sewing kernel factors
through the one-particle space.  The \emph{one-particle sewing
theorem} (Theorem~\ref{thm:heisenberg-one-particle-sewing}):
\[
Z_g(\cH_k)
\;=\;
\det\nolimits_{\mathrm{Fred}}
\bigl(1- S_g^{(1)}\bigr)^{-k},
\]
where $S_g^{(1)}$ is the genus-$g$ one-particle Bergman sewing
operator on the one-particle Hilbert space
$\cH_k^{(1)} = \bigoplus_{n \ge 1}
\mathbb{C}\cdot\alpha_{-n}|0\rangle$.
The determinant is a Fredholm determinant: it converges because
$S_g^{(1)}$ is trace-class (the Bergman kernel
is smooth away from the diagonal and square-integrable).
The formula $Z_g=\det(1-S)^{-k}$ is the exact partition
function at every genus, not a perturbative approximation: each
of the $k$ free bosons contributes $\det(1-S)^{-1}$.  Its
logarithm is
\[
\log Z_g(\cH_k)
\;=\;
-k\,\operatorname{tr}\log(1-S_g^{(1)})
\;=\;
k\sum_{n\ge 1}\frac{1}{n}\,
\operatorname{tr}(S_g^{(1)})^n.
\]
At genus~$1$: $\operatorname{tr}K_1^{(1)}=\log|\eta(\tau)|^{-2}$,
recovering $Z_1(\cH_k)=|\eta(\tau)|^{-k}$.  The one-particle
reduction collapses the infinite-dimensional sewing problem to
a determinant of a single integral operator.

\subsubsection*{The $\hat A$-genus generating function}

On the proved uniform-weight lane
(Definition~\ref{def:scalar-lane}), the scalar-level genus expansion of
the MC element gives the free energies
$F_g(\cA)=\kappa(\cA)\cdot\lambda_g$, where $\lambda_g$ is the
universal coefficient at genus~$g$.  For arbitrary multi-weight
families, only the genus-$1$ identity $F_1=\kappa/24$ is
unconditional.
The generating function assembles into the $\hat A$-genus:
\begin{equation}\label{eq:Ahat-preface-supp}
\sum_{g\ge 1}F_g(\cA)\,x^{2g}
\;=\;
\kappa(\cA)\cdot
\bigl(\hat A(ix)-1\bigr),
\end{equation}
where
$\hat A(x)=x/2\cdot(\sinh(x/2))^{-1}
=1-x^2/24+7x^4/5760-\cdots$
is the Hirzebruch $\hat A$-genus, evaluated at imaginary argument.
The $\hat A$-function appears because the universal coefficients
$\lambda_g=\tfrac{2^{2g-1}-1}{2^{2g-1}}\cdot\tfrac{|B_{2g}|}{(2g)!}$
are Faber--Pandharipande $\lambda_g$-integrals
\cite{FP03}, and $\hat A(x)=(x/2)/\sinh(x/2)$ is their
generating function.

The gravitational interpretation:
equation~\eqref{eq:Ahat-preface-supp} states that, on the proved
scalar lane, the perturbative scalar partition function of the
HT theory is an $\hat A$-genus weighted by the gravitational
central charge $\kappa(\cA)$.  The $\hat A$-genus is the index
density of the Dirac operator; the scalar-lane genus expansion
of $\Theta_\cA$ produces it as a formal consequence of $D^2=0$
on the moduli space of
curves.  The connection between the genus expansion and the
$\hat A$-genus is structural, not coincidental: both are
controlled by the Bernoulli numbers because both arise from the
Todd class of the Hodge bundle via the Grothendieck--Riemann--Roch
theorem applied to the universal curve
$\pi\colon\cC_g\to\overline{\cM}_g$.

\subsubsection*{The sewing envelope}

The perturbative finiteness criterion leads to the notion of
the \emph{sewing envelope} $\cA^{\mathrm{sew}}$: the Hausdorff
completion of~$\cA$ with respect to the family of seminorms
\[
p_{g,q}(a)
\;:=\;
\sup_{b\in\cA,\,\|b\|\le 1}
\bigl|K_g(a,b;q)\bigr|
\]
indexed by genus~$g$ and modular parameter~$q$.  The sewing
envelope carries the analytic continuation of the MC element:
$\Theta_\cA^{\mathrm{an}}\in\MC(\gAmod\,\hat\otimes\,
\cA^{\mathrm{sew}})$.  The algebraic MC element (a formal power
series) converges in $\cA^{\mathrm{sew}}$ whenever the HS
condition holds.  The platonic chain
\[
\cA_{\mathrm{alg}}
\;\subset\;
\cA^{\mathrm{sew}}
\;\rightsquigarrow\;
\cF_\cA\in\mathrm{IndHilb}
\;\rightsquigarrow\;
Q_\bullet^{\mathrm{an}}(\cA)
\]
connects the bare algebraic skeleton to the IndHilb-valued
factorisation homology (in the sense of Moriwaki) and then to
the analytic coderived shadows.  The HS-sewing theorem guarantees
that the first arrow (algebraic to sewing envelope) is an
inclusion for the entire standard landscape of universal algebras
(Remark~\ref{rem:thqg-I-admissible} discusses the admissible-level
simple quotient case).

\subsubsection*{The analytic bar coalgebra}

The HS-sewing criterion leads naturally to an analytic refinement
of the bar construction.  Define the \emph{analytic bar coalgebra}
$\barB^{\mathrm{an}}(\cA)$ as the closure of the algebraic bar
complex $\barB(\cA)$ in the graph-norm topology:
\[
\|\xi\|_{\mathrm{graph}}
\;:=\;
\sum_{n\ge 1}\sum_{a_1,\dots,a_n}
|\xi(sa_1|\cdots|sa_n)|^2\cdot
\prod_i q^{h_i},
\]
where the sum ranges over basis elements of conformal weight
$h_1,\dots,h_n$ and $q=e^{2\pi i\tau}$ is the sewing parameter.
The graph norm makes the bar differential continuous: the OPE
estimate $\mathrm{wt}\,\mu_r(a_1,\dots,a_r)\le
\sum_i\mathrm{wt}(a_i)-(r-1)$ controls the weight decrease under
contraction, and the polynomial OPE growth ensures that the
graph-norm of $d_{\barB}\xi$ is bounded by the graph-norm
of~$\xi$.

The analytic bar coalgebra $\barB^{\mathrm{an}}(\cA)$ is the
natural domain for the genus expansion: the MC element
$\Theta_\cA$ converges in $\barB^{\mathrm{an}}$ precisely when
the HS-sewing condition holds.  The algebraic bar
$\barB(\cA)\subset\barB^{\mathrm{an}}(\cA)$ is a dense
sub-coalgebra; the analytic completion fills in the missing
limits of convergent sequences of graph sums.

\subsubsection*{No UV renormalisation}

The FM compactification resolves all ultraviolet divergences.
When $n$ points collide on a curve~$X$, the
collision is resolved by the iterated blowup
$\overline{\FM}_n(X)\to X^n$: the boundary face $D_S$ records
the collision of the subset $S\subset\{1,\dots,n\}$ as a
smooth codimension-one stratum, not as a singularity.  The
bar differential extracts residues along these smooth boundary
faces; no regularisation is needed because the residue integral
converges.

At higher genus, the moduli space $\overline{\cM}_{g,n}$ is
compact, and the graph-sum integrals are integrals of smooth
forms over compact manifolds.  The only possible divergence is
infrared: the large-modulus regime $\tau\to i\infty$ where
the curve degenerates.  The HS-sewing criterion is precisely the
control of this infrared divergence.  The absence of UV
divergence is geometric (compactness of $\overline{\FM}_n$); the
absence of IR divergence is analytic (HS-sewing).  Together, they
give perturbative finiteness at every genus and every number of
insertions.


% ====================================================================
\subsection*{Gravitational complexity: the $G/L/C/M$ classification}
% ====================================================================

\subsubsection*{Shadow depth as gravitational complexity}

The \emph{shadow depth} $r_{\max}(\cA)$ of a chiral algebra is
the smallest integer~$r$ such that the shadow obstruction tower
terminates: $\Theta_\cA^{\le r}=\Theta_\cA$.  If no such $r$
exists, $r_{\max}=\infty$.  The shadow depth classifies the
complexity of the modular homotopy type, and, by the holographic
dictionary, the complexity of the bulk gravitational theory.

Four classes emerge:

\medskip\noindent
\textbf{Class G} (Gaussian), $r_{\max}=2$.\enspace
The MC element is purely scalar in arity:
$\Theta_\cA=\eta\otimes\Gamma_\cA$ for some genus coefficient
$\Gamma_\cA$.  All higher shadows vanish: $\mathfrak C=0$,
$\mathfrak Q=0$, $\mathrm{Sh}_r=0$ for $r\ge 3$.  The bulk
theory is free (Gaussian): no interaction vertices beyond the
propagator.  On the proved scalar lane this further reduces to
$\Theta_\cA=\kappa\cdot\eta\otimes\Lambda$.

\emph{Representatives}: Heisenberg $\cH_k$, free fermions
$\psi\bar\psi$, and every abelian chiral algebra.

The gravitational content of a class~G theory is entirely
captured by the scalar sector.  On the proved scalar lane, the
genus-$g$ partition function is the $\hat A$-genus weighted
by~$\kappa$; more generally, the Gaussian statement is that no
higher-arity interactions appear.  The shadow connection is flat
with constant curvature, and the modular homotopy type is
formal: no Massey products, no higher obstructions, no contact
interactions.

\medskip\noindent
\textbf{Class L} (Lie/tree), $r_{\max}=3$.\enspace
The cubic shadow $\mathfrak C(\cA)\neq 0$: the Lie bracket
contributes a nontrivial tree-level three-point vertex.  The
quartic shadow vanishes:
$\mathfrak Q(\cA)=0$ by cubic gauge triviality
(Theorem~\ref{thm:cubic-gauge-triviality}), which applies when
$H^1(F^3\gAmod/F^4\gAmod,d_2)=0$.  The bulk theory has
tree-level cubic interactions but no irreducible quartic contact
term.

\emph{Representatives}: affine Kac--Moody
$\widehat{\mathfrak g}_k$ at every non-critical level.  The
universal vertex algebra $V_k(\mathfrak g)$ belongs to class~L
for every simple~$\mathfrak g$ and every
$k\neq-h^\vee$.  Lattice vertex algebras $V_\Lambda$ with
non-degenerate lattice are class~L when the lattice has rank
$\ge 2$ (the rank-$1$ case is Heisenberg, class~G).

The shadow depth $r_{\max}=3$ means the gravitational
theory is ``Lie-type'': its interaction structure is governed
by the Lie algebra (the first-order pole of the OPE), and the
bulk physics is tree-level exact.  The modular homotopy type
has a single nontrivial Massey product (the cubic shadow) but no
higher-order obstructions.

\medskip\noindent
\textbf{Class C} (contact/quartic), $r_{\max}=4$.\enspace
The quartic resonance class $\mathfrak Q(\cA)\neq 0$: there is
an irreducible four-point contact interaction in the bulk.  The
quintic obstruction vanishes.  The tower terminates at the first
genuinely nonlinear shadow beyond the Lie structure.

\emph{Representatives}: $\beta\gamma$ and $bc$ systems at any
value of the conformal dimension.  Both have two generators and the
same stratum separation mechanism that kills the quintic obstruction;
shadow class depends on generator count, not statistics.
The quartic contact invariant of $\beta\gamma$
is
\[
\mathfrak Q^{\mathrm{contact}}_{\beta\gamma}
\;=\;0
\]
by the vanishing theorem
(Corollary~\ref{cor:nms-betagamma-mu-vanishing}), but the full
quartic shadow including the planted-forest correction is
nonzero: the arity-$4$ log-FM residue produces a nontrivial
$\mathrm{Sh}_4$ even when the contact piece vanishes.  The
distinction between $\mathfrak Q^{\mathrm{contact}}=0$ and
$\mathrm{Sh}_4\neq 0$ is characteristic of class~C.

The gravitational content of a class~C theory includes a
single quartic bulk vertex, arising from the codimension-two
boundary of $\overline{\FM}_4(X)$, the first stratum requiring
the log-FM compactification of Mok.  All
higher vertices ($r\ge 5$) are determined by the lower ones
through the MC recursion.

\medskip\noindent
\textbf{Class M} (mixed/infinite), $r_{\max}=\infty$.\enspace
The shadow obstruction tower does not terminate: there exist irreducible
obstruction classes $o^{(r)}(\cA)\neq 0$ at arbitrarily high
arities.  The bulk theory has infinitely many independent
interaction vertices.

\emph{Representatives}: Virasoro $\mathrm{Vir}_c$ at all
$c\neq 0$, $\mathcal W_N$ for $N\ge 2$, non-principal
$\mathcal W$-algebras.  The prototype is
Virasoro: $o^{(5)}_{\mathrm{Vir}}\neq 0$ forces the tower
beyond arity~$4$, and the infinite depth is forced by the
conformal anomaly.

The gravitational content of a class~M theory is the full
shadow obstruction tower: an infinite sequence of irreducible
bulk vertices that cannot be truncated.  The infinite tower
of higher-derivative corrections is the algebraic counterpart
of non-renormalisability.

\subsubsection*{The operadic complexity theorem}

The shadow depth $r_{\max}(\cA)$ should equal the $A_\infty$-depth
of the transferred structure on bar cohomology, and the
$L_\infty$-formality level of the modular convolution algebra.
The precise statement
(Theorem~\ref{thm:operadic-complexity}):
\[
r_{\max}(\cA)
\;=\;
\inf\{n\ge 2\mid m_k=0\;\text{for all}\;k>n\}
\;=\;
\inf\{n\ge 2\mid\ell_k^{(0)}=0\;\text{for all}\;k>n\}.
\]
The first equality identifies shadow depth with $A_\infty$-depth:
the shadow obstruction tower terminates at arity~$r$ iff the minimal
$A_\infty$-structure on $H^*(\barB(\cA))$ has $m_k=0$ for
$k>r$.  The second identifies $A_\infty$-depth with
$L_\infty$-formality level at genus~$0$.  Both identifications
are proved at arities $2$, $3$, $4$
(Proposition~\ref{prop:shadow-formality-low-arity}): the shadow
tower agrees with the $L_\infty$-formality obstruction tower at
low arities.  The conjecture extends this to all arities.

The gravitational consequence: the $G/L/C/M$ classification is
intrinsic: it depends only on the homotopy type of the bar
complex, not on the presentation of~$\cA$.

\subsubsection*{Shadow depth versus Koszulness}

Shadow depth does \emph{not} characterise Koszulness.  Both
finite and infinite shadow-depth algebras can be Koszul:
\begin{center}
\renewcommand{\arraystretch}{1.15}
\begin{tabular}{lcc}
& Koszul & shadow depth \\
\hline
$\cH_k$ & yes & $2$ (class G) \\
$\widehat{\mathfrak g}_k$ & yes & $3$ (class L) \\
$\beta\gamma$ & yes & $4$ (class C) \\
$\mathrm{Vir}_c$ & yes & $\infty$ (class M)
\end{tabular}
\end{center}
Shadow depth classifies the \emph{complexity} of Koszul
algebras (the depth of the nonlinear modular data within the
Koszul locus), not whether an algebra belongs to the Koszul
locus at all.  The twelve characterizations of Koszulness
(Theorem~\ref{thm:koszul-equivalences-meta}: ten unconditional
equivalences, one conditional on perfectness, one one-directional)
are orthogonal to shadow depth.

\subsubsection*{The shadow depth ladder}

Shadow depth increases under natural operations:
\begin{itemize}
\item \emph{DS reduction increases depth}: for
  $\widehat{\mathfrak{sl}}_N\to\mathcal W_N$, the depth jumps
  from~$3$ to~$\infty$.  The BRST functor creates irreducible
  contact vertices from the ghost sector.
\item \emph{Tensor product preserves depth}: for
  $\cA\otimes\cA'$ with vanishing mixed OPE,
  $r_{\max}(\cA\otimes\cA')
  =\max(r_{\max}(\cA),r_{\max}(\cA'))$.
  Shadows factorise by Proposition~\ref{prop:independent-sum-factorization}:
  $\kappa$ is additive, $\Delta$ multiplicative, $R_4^{\mathrm{mod}}$
  additive.
\item \emph{Completion preserves depth}: for the inverse limit
  $\mathcal W_\infty=\varprojlim_N\mathcal W_N$, each finite
  stage has $r_{\max}=\infty$, and the limit inherits
  $r_{\max}=\infty$.  The completion entropy $h_K$ refines the
  infinite-depth class into a continuous invariant.
\end{itemize}

The shadow depth ladder for the standard landscape:
\[
2\;(\cH)\;<\;3\;(\widehat{\mathfrak g})\;<\;4\;(\beta\gamma)
\;<\;\infty\;(\mathrm{Vir},\;\mathcal W_N,\;\mathcal W_\infty).
\]
Each step corresponds to a qualitative change in gravitational
complexity: free $\to$ tree-level $\to$ one-loop contact
$\to$ fully interacting.

\subsubsection*{Completion entropy as gravitational complexity
measure}

Within class~M, the completion entropy
$h_K(\cA)=\log(\rho_\cA^{-1})$ provides a continuous refinement.
The entropy depends only on the vacuum character $\chi_\cA(q)$
through the primitive generating series $G_\cA(t)
=\frac{1-t}{t}(\chi_\cA(t)-1)$.  It measures the exponential
growth rate of the reduced-weight windows $K_q(\cA)$: a higher
entropy means a more complex completion problem, hence a more
complex gravitational theory.

The entropy ladder for $\mathcal W_N$:
\begin{center}
\renewcommand{\arraystretch}{1.15}
\begin{tabular}{lccccl}
$N$ & $2$ & $3$ & $4$ & $\cdots$ & $\infty$ \\
\hline
$h_K$ & $0.567$ & $0.772$ & $0.835$ & $\cdots$ & $0.872$
\end{tabular}
\end{center}
The limit $h_K(\mathcal W_\infty)\approx 0.872$ is controlled
by the MacMahon plane-partition generating function.  The entropy
is monotonically increasing in~$N$: higher-spin gravity is
monotonically more complex than lower-spin gravity, as measured by
the character-level completion kinematics.

For Virasoro, $h_K(\mathrm{Vir}_c)=\log(1)\approx 0.567$
at all~$c$ (the vacuum character of $\mathrm{Vir}_c$ has
partition-function growth independent of~$c$).  The entropy does
not see the central charge; it sees only the conformal-weight
spectrum.


% ====================================================================
\subsection*{$S$-duality, complementarity, and $c=26$}
% ====================================================================

\subsubsection*{Verdier involution as $S$-duality}

The Koszul involution $\cA\mapsto\cA^!$ exchanges boundary and
bulk:
\[
\mathcal H(T)^!
\;=\;
\bigl(\cA^!,\;(\cA^!)^!,\;\cC^!,\;r^!(z),\;
\Theta_{\cA^!},\;\nabla^{\mathrm{hol}}_{\cA^!}\bigr).
\]
For a Koszul algebra, $(\cA^!)^!\simeq\cA$: the involution has
period~$2$.  The boundary of~$T$ is the bulk of~$T^!$ and
vice versa.  In the physical language, this is $S$-duality: the
strong-coupling regime of one theory is the weak-coupling regime
of its dual.

At the level of the MC element:
$\Theta_{\cA^!}=D_{\cA^!}-d_0$, the bar-intrinsic construction
applied to the dual algebra.  The complementarity theorem
(Theorem~C) gives the splitting
\[
R\Gamma(\overline{\cM}_g,\,\cZ_\cA)
\;\simeq\;
Q_g(\cA)\;\oplus\;Q_g(\cA^!),
\]
with $\kappa(\cA)+\kappa(\cA^!)=0$ (for Kac--Moody and
free-field algebras).  The splitting is the
$S$-duality decomposition of the genus-$g$ conformal block:
the boundary contribution~$Q_g(\cA)$ and the bulk
contribution~$Q_g(\cA^!)$ together exhaust the center local
system on~$\overline{\cM}_g$.

\subsubsection*{Lagrangian geometry}

The complementarity pairing upgrades to shifted-symplectic
geometry.  Let
$\omega^{\mathrm{symp}}_{g,n}\in\Omega^2(\cZ_{\cA,g,n},2-2g)$
be the $(-1)$-shifted symplectic form on the center local system
at genus~$g$.  The complementarity splitting
$\cZ_\cA=Q_g(\cA)\oplus Q_g(\cA^!)$ is a \emph{Lagrangian
decomposition}: each summand is Lagrangian with respect to
$\omega^{\mathrm{symp}}$:
\[
\omega^{\mathrm{symp}}\big|_{Q_g(\cA)}=0,
\qquad
\omega^{\mathrm{symp}}\big|_{Q_g(\cA^!)}=0.
\]
The pairing between the two summands is the residue of
$\omega^{\mathrm{symp}}$ along the intersection; this is
the complementarity pairing~$C$.  The Lagrangian condition
is the algebraic statement $\kappa(\cA)+\kappa(\cA^!)=0$
at the scalar level (for Kac--Moody and free-field algebras;
for Virasoro, $\kappa+\kappa^!=13$); at higher arities it
imposes constraints on the shadow obstruction tower of the pair
$(\cA,\cA^!)$.

The gravitational interpretation: the phase space of
3d gravity is the symplectic manifold
$T^*\cM_{\mathrm{flat}}(G,\Sigma_g)$, the cotangent bundle
of the moduli of flat connections.  The Lagrangian decomposition
$Q_g(\cA)\oplus Q_g(\cA^!)$ is the polarisation of this phase
space into position and momentum: the boundary degrees of freedom
and the bulk degrees of freedom.  The complementarity theorem
states that the total Hilbert space at genus~$g$ splits into
boundary and bulk sectors, with the symplectic pairing between
them.

\subsubsection*{The Virasoro self-dual locus}

For the Virasoro family $\mathrm{Vir}_c$:
$\mathrm{Vir}_c^!=\mathrm{Vir}_{26-c}$.  The Koszul involution
acts by $c\mapsto 26-c$.  Self-duality occurs at the fixed point:
\[
c=26-c
\qquad\Longrightarrow\qquad
\boxed{\;c=13\;}
\]
At $c=13$: $\mathrm{Vir}_{13}^!=\mathrm{Vir}_{13}$.  The
boundary and bulk algebras coincide.  The complementarity
splitting degenerates: $Q_g(\cA)=Q_g(\cA^!)$, and the Lagrangian
subspaces merge into a single isotropic subspace.  The
$S$-duality group reduces from $\Z/2$ to the identity.  The
shadow connection has $\kappa(\mathrm{Vir}_{13})=13/2$
and $\kappa(\mathrm{Vir}_{13}^!)=\kappa(\mathrm{Vir}_{13})=13/2$
\textup{(}self-dual\textup{)}, with
$\kappa+\kappa'=13$.

At $c=26$: the dual curvature vanishes:
\[
\kappa(\mathrm{Vir}_{26}^!)
\;=\;\kappa(\mathrm{Vir}_0)
\;=\;\frac{0}{2}\;=\;0.
\]
The Koszul dual $\mathrm{Vir}_0$ is \emph{uncurved}: its bar
complex has $d^2=0$ strictly, with no curvature correction.
The holographic datum at $c=26$:
\begin{align*}
\mathcal H(\mathrm{Vir}_{26})
&\;=\;
\Bigl(
\mathrm{Vir}_{26},\;
\mathrm{Vir}_0,\;
\cC_{26},\;r_{26}(z),\;\Theta_{26},\;\nabla^{\mathrm{hol}}_{26}
\Bigr),\\[4pt]
\kappa_{26}&=13,
\qquad
\kappa_0=0,
\qquad
\mathfrak Q^{\mathrm{contact}}_{26}
=\frac{10}{26\cdot 152}=\frac{5}{1976}\,.
\end{align*}
The critical dimension of the bosonic string emerges as the
vanishing of the dual curvature: $c=26$ is the unique value
where the bulk theory ($\mathrm{Vir}_0$) is uncurved.  The
boundary theory retains curvature $\kappa=13$; it is the
\emph{bulk} that becomes free at $c=26$.  This is the
Koszul-theoretic content of the critical dimension.

\subsubsection*{The $c=13$ versus $c=26$ dichotomy}

The Virasoro family has two distinguished central charges:
\begin{itemize}
\item $c=13$: self-duality.
  $\mathrm{Vir}_{13}\cong\mathrm{Vir}_{13}^!$.  The Koszul
  involution is trivial.  The complementarity pairing becomes
  an inner product.  The shadow connection is
  self-conjugate: $\nabla^{\mathrm{hol}}_{13}
  =(\nabla^{\mathrm{hol}}_{13})^!$.
\item $c=26$: uncurved dual.
  $\mathrm{Vir}_0$ has $\kappa=0$ and strictly nilpotent bar
  differential.  The bulk theory is free.  The genus expansion
  of the dual is trivial beyond genus~$0$: $F_g(\mathrm{Vir}_0)
  =0$ for $g\ge 1$.
\end{itemize}
These are different conditions, and they do not coincide.
Self-duality ($c=13$) is a symmetry condition; uncurved dual
($c=26$) is a finiteness condition.  The bosonic string sits at
the finiteness point, not the symmetry point.

For the $\mathcal W_N$ family at $N=3$: the Koszul dual is
$\mathcal W_3{}_{50-c}$.  Self-duality at $c=25$.  Uncurved dual at
$c=50$ (where $\kappa(\mathcal W_3{}_{50})=0$).  The pattern
generalises: each $\mathcal W_N$ family has a self-dual central
charge $c_{\mathrm{sd}}(N)$ and a critical central charge
$c_{\mathrm{crit}}(N)$, with
$c_{\mathrm{sd}}\neq c_{\mathrm{crit}}$ in general.

\subsubsection*{The Feigin--Frenkel involution}

For affine Kac--Moody $\widehat{\mathfrak g}_k$, the Koszul
involution acts on the modular characteristic by the
Feigin--Frenkel level shift $k\mapsto-k-2h^\vee$:
the Koszul dual $(\widehat{\mathfrak g}_k)^!$
has $\kappa((\widehat{\mathfrak g}_k)^!)
= \dim\mathfrak g\cdot(-k-h^\vee)/(2h^\vee)$,
matching $\kappa(\widehat{\mathfrak g}_{-k-2h^\vee})$,
but $(\widehat{\mathfrak g}_k)^! \neq
\widehat{\mathfrak g}_{-k-2h^\vee}$ as chiral algebras
(AP33: the Koszul dual is the linear dual of bar cohomology,
not the affine algebra at the shifted level).
The critical level $k=-h^\vee$ is where the
Sugawara construction is undefined ($\kappa = 0$) and the
center becomes the algebra of $\mathfrak g$-opers
(Feigin--Frenkel).  The Koszul involution on levels
$k\mapsto -k-2h^\vee$ has unique fixed point $k=-h^\vee$,
which is singular, not smooth: the affine Koszul
involution has no smooth fixed point at generic level.

The complementarity:
$\kappa(\widehat{\mathfrak g}_k)
+\kappa(\widehat{\mathfrak g}_{-k-2h^\vee})=0$.
Explicitly for $\mathfrak{sl}_2$:
\[
\kappa_k=\frac{3(k+2)}{4},
\qquad
\kappa_{-k-4}=\frac{3(-k-4+2)}{4}
=\frac{3(-k-2)}{4}=-\frac{3(k+2)}{4}\,,
\]
and $\kappa_k+\kappa_{-k-4}=3(k+2)/4-3(k+2)/4
=0$.  The curvatures sum to zero: this is the
complementarity identity $\kappa+\kappa^!=0$
for the affine Kac--Moody family (for Virasoro,
$\kappa+\kappa^!=13$; AP24).

\subsubsection*{The depth-zero resonance shadow}

The Koszul dual $\mathrm{Vir}_{26-c}$ is the
\emph{depth-zero resonance shadow}: it is the image of the
finite-dimensional resonance truncation
$R_{\mathrm{Vir}}$ at depth zero in the shadow obstruction tower.  It is
\emph{not} the final Koszul dual in the non-perturbative sense;
the genuine $\mathcal W_\infty$-type dual requires the full
resonance-filtered completion
(Theorem~\ref{thm:platonic-completion}).  The distinction is
between the same-family shadow (a Virasoro algebra at shifted
central charge) and the true categorical dual (an object in the
completed category).

At $c=26$: the depth-zero resonance shadow is $\mathrm{Vir}_0$,
which is the trivial conformal field theory (the stress tensor
acts as zero on the vacuum module).  The string-theoretic
statement ``the ghost central charge cancels the matter central
charge'' is the Koszul-theoretic statement ``the dual curvature
vanishes at $c=26$.''

\subsubsection*{Complementarity at higher genus}

At genus $g\ge 2$, the complementarity splitting acquires
corrections from the planted-forest strata:
\begin{align*}
R\Gamma(\overline{\cM}_g,\cZ_\cA)
&\;\simeq\;
Q_g(\cA)\;\oplus\;Q_g(\cA^!)\;\oplus\;Q_g^{\mathrm{mix}}(\cA),\\[4pt]
Q_g^{\mathrm{mix}}(\cA)
&\;=\;
\bigoplus_{\Gamma\in\mathsf{Gr}^{\mathrm{st,mix}}_{g}}
\frac{1}{|\operatorname{Aut}(\Gamma)|}
\;W_\Gamma^{\mathrm{mix}},
\end{align*}
where $\mathsf{Gr}^{\mathrm{st,mix}}_g$ is the set of stable
graphs with mixed vertex decorations: some vertices decorated
by~$\cA$, others by~$\cA^!$.  The mixed term is controlled by
the complementarity pairing~$C$ and vanishes at genus~$1$
(where the only stable graph is the self-loop, which carries
a single vertex).  At genus~$2$, the mixed separating-node
graph $\Theta^{(2)}_{\mathrm{sep}}$ has two vertices
(genera~$1$ and~$1$, joined by one edge), and the mixed
decoration assigns $\cA$ to one vertex and~$\cA^!$ to the other.
Its contribution is $\kappa(\cA)\cdot\kappa(\cA^!)\cdot
C(\alpha,\alpha)$, a product of curvatures weighted by the
complementarity pairing.

For Kac--Moody and free-field algebras, the condition
$\kappa(\cA)+\kappa(\cA^!)=0$ ensures that the
mixed contribution simplifies: the product
$\kappa(\cA)\kappa(\cA^!)=-\kappa(\cA)^2$ is negative definite.
(For Virasoro, $\kappa+\kappa^!=13\ne 0$, so the mixed term
$\kappa\kappa^!=\kappa(13-\kappa)$ is not sign-definite.)
The higher-genus complementarity is the statement that these
mixed terms, together with the direct boundary and bulk terms,
exhaust the moduli-space cohomology.


% ====================================================================
\subsection*{The gravitational Yangian}
% ====================================================================

\subsubsection*{Collision residue and the $r$-matrix}

The genus-$0$, arity-$2$ component of $\Theta_\cA$, evaluated
at two points $z_1,z_2\in\C$, defines a kernel
$\Theta_\cA^{(0,2)}(z_1,z_2)
\in\cA\otimes\cA\otimes\Omega^1_{\mathrm{log}}$.
The \emph{collision $r$-matrix} is its residue along the
diagonal:
\[
r(z)
\;:=\;
\Res_{z_1\to z_2}\Theta_\cA^{(0,2)}(z_1,z_2)
\;=\;
\Res^{\mathrm{coll}}_{0,2}(\Theta_\cA).
\]
This is a meromorphic function of $z=z_1-z_2$ valued in
$\cA\otimes\cA$: the collision kernel of the OPE, promoted to
an element of the convolution algebra.

For the Heisenberg algebra: $r(z)=k/z$, the constant Casimir
divided by the separation.  For affine Kac--Moody:
$r(z)=\Omega/z$, the full Casimir element.  For Virasoro:
$r(z)=c/(2z^3)+2T/z$, with a third-order pole reflecting the
fourth-order pole $T(z)T(w)\sim c/2(z-w)^{-4}+\cdots$ in the OPE
(the $d\log$ extraction absorbs one pole order).%
\footnote{This is the \emph{pre-dualisation form}
$r^{\mathrm{pre}}(z)$ of
Remark~\textup{\ref{rem:three-r-matrices}}\textup{(c)}, carrying
collision-residue poles and valued in $\cA\otimes\cA$.
After Koszul dualisation, the same data gives
$r^{\mathrm{coll}}(z) \in \cA^!\otimes\cA^!$.  The
\emph{$E_1$ scalar shadow} $r^{\mathrm{sc}}(z) = c/(2z)$
appearing in the archetype table
\textup{(}Example~\textup{\ref{ex:e1-shadow-archetypes})}
is the vacuum projection to a single scalar channel.}

The $r$-matrix is the genus-$0$ binary shadow
$r(z)=\mathrm{Sh}_{0,2}(\Theta_\cA)(z)$.  It is the
lowest-arity, lowest-genus projection of the MC element, and
it controls the tree-level scattering of the bulk theory.

\subsubsection*{Arnold implies Yang--Baxter}

The classical Yang--Baxter equation for~$r(z)$ is a direct
consequence of the Arnold relation.  The argument proceeds in
three steps.

\emph{Step 1.}\enspace
Consider $\Theta_\cA^{(0,3)}(z_1,z_2,z_3)$, the genus-$0$,
arity-$3$ component of the MC element.  The MC equation at
genus~$0$ and arity~$3$ is
\begin{equation}\label{eq:mc-g0-a3-preface-supp}
d_0\,\Theta_\cA^{(0,3)}
+\tfrac{1}{2}[\Theta_\cA^{(0,2)},
\Theta_\cA^{(0,2)}]
\;=\;0.
\end{equation}
The bracket is the genus-$0$ bracket on $\gAmod$, which is the
tree-level graph-sum bracket.

\emph{Step 2.}\enspace
The bracket $[\Theta^{(0,2)},\Theta^{(0,2)}]$ evaluated at
$(z_1,z_2,z_3)$ sums over the three trees on three leaves:
\begin{align*}
\tfrac{1}{2}[\Theta^{(0,2)},\Theta^{(0,2)}]_{123}
&\;=\;
r_{12}\star r_{13}
+r_{12}\star r_{23}
+r_{13}\star r_{23}\\[4pt]
&\;=\;
[r_{12},r_{13}]\,\eta_{12}\wedge\eta_{13}
+[r_{12},r_{23}]\,\eta_{12}\wedge\eta_{23}\\
&\qquad\qquad
+[r_{13},r_{23}]\,\eta_{13}\wedge\eta_{23},
\end{align*}
where the $\star$-product inserts the $r$-matrix at each internal
edge of the tree and contracts by the chiral pairing.

\emph{Step 3.}\enspace
Apply the Arnold relation:
$\eta_{12}\wedge\eta_{23}
+\eta_{23}\wedge\eta_{31}
+\eta_{31}\wedge\eta_{12}=0$.
Substituting $\eta_{31}=-\eta_{13}$:
\[
\eta_{12}\wedge\eta_{23}
-\eta_{23}\wedge\eta_{13}
-\eta_{13}\wedge\eta_{12}
=0.
\]
The MC equation~\eqref{eq:mc-g0-a3-preface-supp}, after extracting
the residue along all three diagonals simultaneously, becomes
\[
[r_{12},r_{13}]+[r_{12},r_{23}]+[r_{13},r_{23}]
\;=\;0.
\]
This is the CYBE.  The Arnold relation on $\overline{\FM}_3(\C)$
is the geometric origin of the Yang--Baxter equation; the CYBE
is the $d^2=0$ identity at genus~$0$ and arity~$3$, viewed
through the collision residue.

\subsubsection*{The dg-shifted Yangian}

The \emph{dg-shifted Yangian} $\Ydg(\cA)$ is the algebra
generated by the modes of the $r$-matrix: formally, the
$\Ran(\C)$-envelope of the OPE algebra of~$r(z)$.  For affine
Kac--Moody $\widehat{\mathfrak g}_k$, $\Ydg(\cA)$ is the
Yangian $Y(\mathfrak g)$ with a dg shift recording the
homological degree.  For Virasoro, $\Ydg(\mathrm{Vir}_c)$ is
the dg-shifted stress-tensor Yangian candidate attached to the
genus-zero binary shadow.  On the generic theorem surface of the
manuscript, the independently proved algebraic line-side model is
$\mathrm{Vir}_c^!\text{-}\mathrm{mod}$; a direct dg-shifted-Yangian
realization is only available on the affine / evaluation-core lane.

The Yangian acts on the line-operator category: the
Kazhdan--Lusztig equivalence
$\cC\simeq Y(\mathfrak g)\text{-mod}^{\mathrm{fd}}$
identifies line operators with finite-dimensional
$Y(\mathfrak g)$-modules in the affine case.  For Virasoro, the
analogous dg-shifted-Yangian statement is a target of the
Yangian-shadow programme, not a theorem proved on the present
generic surface.

The modular extension: define $\mathcal{Y}^{\mathrm{dg,mod}}(\cA)$
as the ambient pronilpotent completion of the modular convolution
dg Lie algebra.  A conjectural modular $R$-matrix
$R^{\mathrm{mod}}_\cA(z;\hbar)$ would be an MC element
in $\mathcal{Y}^{\mathrm{dg,mod}}(\cA)$:
\[
R^{\mathrm{mod}}_\cA(z;\hbar)
\;=\;
r(z)+\hbar\,r_1(z)+\hbar^2\,r_2(z)+\cdots\,,
\]
where $r_g(z)$ is the genus-$g$ correction to the $r$-matrix.
On that conjectural surface, the quantum Yang--Baxter equation
$R_{12}R_{13}R_{23}=R_{23}R_{13}R_{12}$ would be the corresponding
MC equation for $R^{\mathrm{mod}}$.  The genus expansion of the
$r$-matrix is intended to parallel the genus expansion of the shadow
package, but this higher-genus modular-kernel identification is not
proved in the present chapter.

\subsubsection*{The gravitational $r$-matrix for Virasoro}

The Virasoro $r$-matrix is
\[
r_c(z)
\;=\;
\frac{c/2}{z^3}+\frac{2T}{z}\,.
\]
The third-order pole $c/(2z^3)$ carries the central extension; the
simple pole $2T/z$ carries the stress-tensor exchange.  The CYBE
for this $r$-matrix:
\begin{align*}
&[r_{12}(z_{12}),r_{13}(z_{13})]
+[r_{12}(z_{12}),r_{23}(z_{23})]
+[r_{13}(z_{13}),r_{23}(z_{23})]\\[4pt]
&\qquad=\;
\frac{c}{z_{12}^3}\Bigl(\frac{2T_3}{z_{13}}
-\frac{2T_3}{z_{23}}\Bigr)
+\frac{4}{z_{12}\,z_{13}}[T_2,T_3]
+\cdots
\;=\;0,
\end{align*}
using $[T(z),T(w)]=(\partial_w T)(w)\cdot\delta(z-w)
+2T(w)\cdot\delta'(z-w)+(c/2)\delta'''(z-w)$.
The cancellation requires all three pole orders of the OPE
simultaneously.

At genus~$1$, the corrected $r$-matrix acquires a modular
correction:
\[
r_c^{(1)}(z;\tau)
\;=\;
r_c(z)+\hbar\cdot\kappa_c\cdot\wp(z;\tau),
\]
where $\wp(z;\tau)$ is the Weierstrass $\wp$-function, the
genus-$1$ propagator.  The quantum YBE at genus~$1$ is the
condition that the $\wp$-correction is compatible with the MC
structure at one loop.  The heat equation for $\wp$ translates
into the flatness of $\nabla^{\mathrm{hol}}_{1,2}$.

\subsubsection*{The $r$-matrix hierarchy and KZ at higher genus}

The collision residue construction at genus~$g$ produces a
family of $r$-matrices $r^{(g)}(z;\Sigma)$ depending on the
Riemann surface $\Sigma\in\overline{\cM}_g$:
\[
r^{(g)}(z;\Sigma)
\;=\;
\Res^{\mathrm{coll}}_{g,2}(\Theta_\cA)(z;\Sigma).
\]
At genus~$0$: $r^{(0)}(z)=r(z)$, the standard $r$-matrix.
At genus~$1$: $r^{(1)}(z;\tau)=r_c^{(1)}(z;\tau)$, the
$\wp$-corrected matrix.  At genus~$g\ge 2$: $r^{(g)}$ depends
on all $3g-3$ moduli of~$\Sigma_g$ through the genus-$g$
Bergman kernel.

The shadow connection at genus~$g$ and arity~$n$ has a KZ-like
binary leading term with $r^{(g)}$ as its structure matrices,
together with separate higher-arity corrections:
\[
\nabla^{\mathrm{hol}}_{g,n}
\;=\;
d\;-\;\sum_{i<j}r^{(g)}_{ij}(z_i-z_j;\Sigma)
\,d\log(z_i-z_j)
\;-\;\sum_{r\ge 3}\mathrm{Sh}_{g,r}(\Theta_\cA).
\]
The higher-arity shadows $\mathrm{Sh}_{g,r}$ for $r\ge 3$ are the
non-KZ corrections: they arise from stable graphs with more than
two external legs and contribute beyond the binary $r$-matrix
term.  The flatness condition for this connection imposes
constraints on $r^{(g)}$ that generalise the CYBE to higher genus.


% ====================================================================
\subsection*{Soft theorems from the shadow obstruction tower}
% ====================================================================

\subsubsection*{Ward identities as shadow projections}

The modular shadow connection
$\nabla^{\mathrm{hol}}_{g,n}=d-\sum_r\mathrm{Sh}_{g,r}
(\Theta_\cA)$ encodes the Ward identities of the bulk theory.
A \emph{soft theorem} is the leading-order Ward identity when one
of the $n$ insertions becomes ``soft,'' when its momentum
(equivalently, conformal weight) goes to zero.

In the chiral framework, the soft limit corresponds to the
collision limit $z_n\to z_{n-1}$: as two insertion points
approach, the OPE extracts a residue, and the leading term of the
residue is the soft factor.  The shadow connection's dependence on
$z_n$ near the collision divisor $D_{\{n-1,n\}}$ gives:
\[
\mathrm{Sh}_{g,n}(\Theta_\cA)\big|_{z_n\to z_{n-1}}
\;\sim\;
\frac{r_{(n-1)n}(z_{n-1}-z_n)}{z_{n-1}-z_n}\cdot
\mathrm{Sh}_{g,n-1}(\Theta_\cA)
\;+\;O(1).
\]
The residue $r_{(n-1)n}/(z_{n-1}-z_n)$ is the soft factor,
and the Ward identity it imposes on the $(n-1)$-point shadow
is the soft theorem.

\subsubsection*{Soft graviton from Virasoro}

For Virasoro, the soft limit of the stress tensor is
\[
\mathrm{Sh}_{0,n+1}(\Theta_c)\big|_{z_{n+1}\to z_n}
\;\sim\;
\Bigl(\frac{h_n}{(z_{n+1}-z_n)^2}
+\frac{\partial_{z_n}}{z_{n+1}-z_n}\Bigr)
\cdot\mathrm{Sh}_{0,n}(\Theta_c),
\]
where $h_n$ is the conformal weight at the $n$-th insertion.
This is the conformal Ward identity: the stress tensor acts as
the $L_{-1}$ (translation) and $L_0$ (dilation) generators in
the soft limit.  In the bulk gravitational language, these
are the BMS supertranslation and superrotation soft theorems.

The connection to the BMS group is structural:
the asymptotic symmetry algebra of 3d gravity at null
infinity is the BMS$_3$ algebra, which is a central extension of
the diffeomorphism group of the circle.  The Virasoro algebra
$\mathrm{Vir}_c$ is the holomorphic part of BMS$_3$.  The soft
graviton theorems (the Ward identities of the BMS$_3$ symmetry)
are the leading-order projections of the shadow connection at
genus~$0$.

At higher genus, the soft theorems acquire corrections from the
period structure.  At genus~$1$:
\[
\mathrm{Sh}_{1,n+1}(\Theta_c)\big|_{z_{n+1}\to z_n}
\;\sim\;
\Bigl(\frac{h_n}{(z_{n+1}-z_n)^2}
+\frac{\partial_{z_n}}{z_{n+1}-z_n}\Bigr)
\cdot\mathrm{Sh}_{1,n}(\Theta_c)
\;+\;\kappa_c\cdot\omega_1(z_n)\cdot
\mathrm{Sh}_{0,n}(\Theta_c),
\]
where $\omega_1(z_n)$ is the normalised abelian differential
evaluated at~$z_n$.  The period correction $\kappa_c\omega_1$
is the genus-$1$ soft graviton theorem: it states that the
stress tensor creates a genus-$1$ tadpole proportional to the
curvature.

\subsubsection*{Subleading soft from quartic shadow}

The quartic contact invariant produces a subleading soft theorem.
The shadow connection at genus~$0$ and arity~$4$, in the limit
$z_4\to z_3$, gives:
\[
\mathrm{Sh}_{0,4}(\Theta_c)\big|_{z_4\to z_3}
\;\sim\;
\frac{r_{34}}{z_{34}}\cdot\mathrm{Sh}_{0,3}(\Theta_c)
\;+\;
\frac{\mathfrak Q_c}{z_{34}}\cdot\delta_4
\;+\;O(1),
\]
where $z_{34}=z_3-z_4$ and $\delta_4$ is the
codimension-$2$ contact correction from the log-FM boundary.
The subleading soft factor $\mathfrak Q_c/z_{34}$ is the
quartic contact shadow $10/[c(5c+22)]$ times the log-FM
correction.  At $c=-22/5$ (the $(2,5)$ minimal model), the
subleading soft factor diverges, a resonance point where the
soft expansion breaks down.

The hierarchy of soft theorems mirrors the shadow obstruction
tower:
\begin{center}
\renewcommand{\arraystretch}{1.15}
\begin{tabular}{lll}
\textbf{Shadow level} & \textbf{Soft theorem}
  & \textbf{Physical content} \\
\hline
$\kappa$ (arity $2$) & leading soft
  & BMS supertranslation \\[2pt]
$\mathfrak C$ (arity $3$) & subsubleading soft
  & stress-tensor three-point coupling \\[2pt]
$\mathfrak Q$ (arity $4$) & contact soft
  & four-graviton contact interaction \\[2pt]
$\mathrm{Sh}_r$ (arity $r$) & order-$(r-2)$ soft
  & $r$-graviton vertex
\end{tabular}
\end{center}
Each level of the shadow obstruction tower produces one order of the
soft expansion.  The infinite depth of the Virasoro tower
(class~M) corresponds to the infinite sequence of subleading
soft theorems in 3d gravity, each one encoding an independent
gravitational interaction.

\subsubsection*{Soft theorems for affine Kac--Moody}

For $\widehat{\mathfrak g}_k$ (class~L, shadow depth~$3$), the
soft hierarchy truncates after the cubic level:
\begin{align*}
\text{leading soft}&:\quad
J^a(z)\sim\frac{f^{ab}{}_c\,J^c(w)}{z-w}
+\frac{k\kappa^{ab}}{(z-w)^2}\,,\\[4pt]
\text{subleading soft}&:\quad
\text{vanishes (cubic gauge triviality)}.
\end{align*}
The subleading soft theorem is trivial because the quartic shadow
vanishes.  The soft expansion for an affine algebra terminates at
leading order: no contact corrections, no higher soft factors.
This reflects the tree-level exactness of the bulk Chern--Simons
theory.

For Virasoro (class~M), the
soft expansion never terminates: every arity contributes an
independent soft factor, and the gravitational $S$-matrix has
infinitely many soft theorems.  For gauge theory (class~L), the
soft expansion terminates after the first nontrivial order: the
Yang--Mills $S$-matrix has only the leading (Weinberg) and
subleading soft theorems, and no higher-order contact soft
factors.

\subsubsection*{Soft theorems and the $\hat A$-genus}

On the proved scalar lane, the generating function for the
scalar soft factors at all genera is controlled by the
$\hat A$-genus.  The genus-$g$ leading soft factor is
$\kappa(\cA)\cdot\lambda_g$, and the generating function
$\sum_{g\ge 1}\lambda_g\,x^{2g}=\hat A(ix)-1$
assembles the soft factors into the $\hat A$-class.  The
$\hat A$-genus is therefore the universal soft factor for the
scalar-level shadow on that lane.  For arbitrary multi-weight
families, only the genus-$1$ soft factor is currently
unconditional.

\subsubsection*{The double-soft limit and the shadow bracket}

The double-soft limit (the simultaneous collision of two
insertions $z_n\to z_{n-1}\to z_{n-2}$) probes the bracket
structure of the shadow connection.  The leading double-soft
factor is
\[
\mathrm{Sh}_{0,n+2}(\Theta_\cA)\big|_{z_{n+2}\to z_{n+1}
\to z_n}
\;\sim\;
\frac{[r_{n(n+1)},r_{n(n+2)}]}{z_{n(n+1)}\cdot z_{n(n+2)}}
\cdot\mathrm{Sh}_{0,n}(\Theta_\cA),
\]
where the commutator $[r,r]$ is the Lie bracket of the
collision residues.  This is a double-soft graviton theorem:
it states that the simultaneous emission of two soft gravitons
is controlled by the Lie algebra structure of the $r$-matrix.
The CYBE guarantees that the double-soft limit is consistent:
the three orderings of the double collision produce the same
result, up to the Arnold relation.

For affine Kac--Moody, the double-soft factor is
$[r,r]=[\Omega/z_1,\Omega/z_2]$, which reduces to
$[\Omega,\Omega]/(z_1z_2)=0$
since the Casimir is central: $[\Omega,\Omega]=0$.
The double-soft graviton theorem is trivial for gauge theories:
there is no simultaneous double-soft emission.  For Virasoro,
$[r_c(z_1),r_c(z_2)]\neq 0$ because the stress tensor does not
commute with itself: $[T(z),T(w)]\neq 0$.  The double-soft
graviton theorem is nontrivial for gravity.

\subsubsection*{Soft theorems and the shadow depth hierarchy}

The number of independent soft theorems equals the shadow depth:
a class~G algebra has one soft theorem (the leading scalar soft
factor); a class~L algebra has two (leading and cubic); a class~C
algebra has three (leading, cubic, and quartic); a class~M algebra
has infinitely many.  The hierarchy is:
\begin{center}
\renewcommand{\arraystretch}{1.15}
\begin{tabular}{lcccc}
\textbf{Class} & G & L & C & M \\
\hline
independent soft theorems & $1$ & $2$ & $3$ & $\infty$ \\
bulk vertices & $0$ & $1$ & $2$ & $\infty$ \\
$r_{\max}$ & $2$ & $3$ & $4$ & $\infty$
\end{tabular}
\end{center}
The number of bulk vertices is $r_{\max}-2$: the first two arities
are the propagator and the two-point function (kinematic data),
and each higher arity contributes one dynamical interaction vertex.
The gravitational complexity is the number of independent dynamical
vertices, which is exactly the number of subleading soft theorems.


% ====================================================================
\subsection*{The modular bootstrap}
% ====================================================================

\subsubsection*{The genus spectral sequence}

The MC equation $D\Theta+\frac{1}{2}[\Theta,\Theta]=0$ admits a
spectral sequence in genus, with $E_1$-page:
\[
E_1^{p,q}
\;=\;
H^q\bigl(\mathfrak{g}_{\cA}^{\mathrm{mod},(p)}\bigr),
\]
where $\mathfrak{g}_{\cA}^{\mathrm{mod},(p)}$ is the genus-$p$ summand of the convolution
algebra.  The $E_1$-page isolates the data by genus:
$p=0$ is tree-level, $p=1$ is one-loop, $p=2$ is genus-$2$.  The
differentials $d_r\colon E_r^{p,q}\to E_r^{p+r,q-r+1}$ are the
obstruction maps: $d_1$ is the genus-$0$ differential, $d_2$ is
the genus-$1$ obstruction, and so on.

The modular bootstrap is the procedure of solving the MC equation
inductively in genus:

\smallskip\noindent
\emph{Step 0} (genus~$0$): solve
$D_0\alpha_0+\frac{1}{2}[\alpha_0,\alpha_0]_0=0$.  The
solution exists by the bar-intrinsic construction:
$\alpha_0=\Theta_\cA^{(0)}$.

\smallskip\noindent
\emph{Step 1} (genus~$1$): solve
$D_0\alpha_1+[\alpha_0,\alpha_1]_0+\hbar\Delta\alpha_0=0$.  The
obstruction is
$[\hbar\Delta\alpha_0]\in E_2^{1,1}$.  For every algebra
in the standard landscape, the obstruction is killed by the
curvature correction $\kappa\cdot\omega_1$: the solution
$\alpha_1=\Theta_\cA^{(1)}$ exists.

\smallskip\noindent
\emph{Step $g$} (genus~$g$): solve
\[
D_0\alpha_g
+\sum_{\substack{g_1+g_2=g\\g_1,g_2\ge 0}}
[\alpha_{g_1},\alpha_{g_2}]_0
+\hbar\,\Delta\alpha_{g-1}
+\text{(clutching terms)}
\;=\;0.
\]
The clutching terms arise from graphs with $b_1\ge 2$ internal
loops; they involve the MC data at lower genera in pairs.  The
obstruction lives in $E_g^{g,1}$ and is killed by the
genus-$g$ component of the Hodge bundle.

The inductive solvability is guaranteed by the bar-intrinsic
construction: $\Theta_\cA=D_\cA-d_0$ is MC at all genera
simultaneously because $D_\cA^2=0$.  The spectral sequence
converges at $E_\infty$ and the totality of solutions assembles
into $\Theta_\cA$.

\subsubsection*{MC recursion as gravitational bootstrap}

The genus spectral sequence is the gravitational bootstrap:
solving the MC equation genus by genus is computing the
perturbative expansion of the gravitational path integral loop
by loop.  The tree-level solution ($g=0$) is the classical
field equation.  The one-loop solution ($g=1$) is the Casimir
energy, the curvature $\kappa(\cA)$.  The genus-$g$ solution
is the $g$-loop contribution to the gravitational partition
function.

The MC recursion
\[
D_0\Theta^{(g)}
\;=\;
-\sum_{\substack{g_1+g_2=g\\g_i<g}}
[\Theta^{(g_1)},\Theta^{(g_2)}]
-\hbar\Delta\Theta^{(g-1)}
-\sum_{\Gamma\in\mathsf{Cl}_g}
\ell_\Gamma(\Theta^{(<g)})
\]
determines $\Theta^{(g)}$ from the lower-genus data, up to
the ambiguity of an element of $\ker D_0$ at genus~$g$.  The
ambiguity is the \emph{modular anomaly}: it measures the freedom
in the gravitational partition function at genus~$g$ that is not
fixed by lower-genus data.  For Virasoro, the ambiguity at
genus~$1$ is a single constant (the choice of regularisation of
the Casimir energy); at genus~$2$ it is a three-dimensional
space corresponding to the three shells
$\Theta^{(2)}_{\mathrm{loop}^2}$,
$\Theta^{(2)}_{\mathrm{sep}}$,
$\Theta^{(2)}_{\mathrm{pf}}$.

The two-dimensional convergence theorem
(MC5, Theorem~\ref{thm:general-hs-sewing}):
the MC recursion converges in both the genus and arity
directions simultaneously.  No UV renormalisation is needed:
the FM compactification resolves all collisions, and the
HS-sewing criterion controls the genus summation.  The
gravitational bootstrap is perturbatively well-defined at all
loop orders.

\subsubsection*{Constraints from MC at genus~$2$}

At genus~$2$, the MC equation imposes nontrivial constraints
on the genus-$1$ data.  Specifically:
\[
D_0\Theta^{(2)}
+[\Theta^{(0)},\Theta^{(2)}]+[\Theta^{(1)},\Theta^{(1)}]
+\hbar\Delta\Theta^{(1)}
+\text{(clutching)}
\;=\;0.
\]
The term $[\Theta^{(1)},\Theta^{(1)}]$ couples the genus-$1$
shadow to itself through the tree-level bracket.  Consistency
at genus~$2$ constrains the genus-$1$ solution: not every
choice of $\alpha_1\in\ker D_0$ at genus~$1$ admits a genus-$2$
extension.  This is the genus-$2$ bootstrap constraint.

For Virasoro at genus~$2$, the three shells contribute:
\begin{align*}
\Theta^{(2)}_{\mathrm{loop}^2}
&\;=\;\frac{c}{2}\cdot\frac{c}{2}\cdot
  (\eta\otimes\Lambda)^{\star 2},\\[4pt]
\Theta^{(2)}_{\mathrm{sep}\circ\mathrm{loop}}
&\;=\;\frac{c}{2}\cdot
  (\eta\otimes\Lambda)\star(\Theta^{(1)}_{\mathrm{Lie}}),\\[4pt]
\Theta^{(2)}_{\mathrm{pf}}
&\;=\;
  \mathrm{Res}_{\mathrm{log-FM}}(\Theta^{(1)}\star\Theta^{(1)}),
\end{align*}
where $\star$ is the convolution product and the planted-forest
term arises from the log-FM residue at the codimension-$2$
boundary of $\overline{\cM}_{2,n}$.  The consistency
of these three shells is a nontrivial check on the MC structure
at genus~$2$.

\subsubsection*{The bootstrap hierarchy for $\mathcal W_N$}

The $\mathcal W_N$ family illustrates the bootstrap at higher rank.
The genus-$1$ Hessian is an $(N-1)\times(N-1)$ matrix
$H^{(1)}_{\mathcal W_N}$, whose eigenvalues control the spectral
discriminant.  The genus-$2$ MC equation constrains the eigenvalues
of $H^{(1)}$ through the quadratic coupling
$[H^{(1)},H^{(1)}]$.  As $N$ increases, the number of independent
genus-$2$ constraints grows quadratically: the bootstrap becomes
increasingly overdetermined.

In the inverse limit $N\to\infty$, the bootstrap hierarchy for
$\mathcal W_\infty$ involves infinitely many generators and
infinitely many genus-$2$ constraints.  The strong
completion-tower theorem (Theorem~\ref{thm:completed-bar-cobar-strong})
guarantees that the completed MC recursion converges: each
finite-window $K_q(\mathcal W_\infty)$ stabilises at
$N=q+1$, and the completed recursion inherits the convergence
from the finite stages.

\subsubsection*{The modular anomaly and the Torelli map}

The ambiguity in the MC recursion (the kernel of $D_0$ at each
genus) is the modular anomaly.  At genus~$g$, the anomaly is an
element of $H^0(\mathfrak{g}_{\cA}^{\mathrm{mod},(g)})$: a genus-$g$ scalar that is not
determined by lower-genus data.  For the Heisenberg algebra
(class~G), the anomaly vanishes at all genera: the scalar
$\kappa$ determines everything.  For Virasoro (class~M), the
anomaly at genus~$g$ is a function on the Torelli space
$\mathcal{T}_g=\overline{\cM}_g\setminus D_{\mathrm{irr}}$, where
$D_{\mathrm{irr}}$ is the irreducible boundary divisor.

The Torelli map $t_g\colon\cM_g\to\cA_g$ to the moduli of
principally polarised abelian varieties (the Jacobian map)
factors the anomaly: the genus-$g$ ambiguity depends only on
the period matrix $\Omega_g$ of the curve through
$t_g$.  This is because the sewing kernel~$K_g$ depends on
$\Sigma_g$ only through its periods (the Bergman kernel is
determined by the period matrix).  The modular anomaly therefore
descends to a function on $\cA_g$.  For $g=1$:
$\cA_1\cong\mathbb H/\mathrm{SL}_2(\Z)$, and the anomaly is
a scalar constant.  For $g=2$: $\cA_2$ is three-dimensional
(Siegel upper half-space modulo $\mathrm{Sp}_4(\Z)$), and the
anomaly has three degrees of freedom, corresponding to the
three shells of the MC element at genus~$2$.

The modular bootstrap progressively eliminates anomaly freedom:
each genus imposes new constraints on the lower-genus data, and
the genus spectral sequence converges to the unique MC element
$\Theta_\cA$.  The bar-intrinsic construction guarantees
existence; the bootstrap shows that the solution is forced by
the $D^2=0$ identity.


% ====================================================================
\subsection*{The critical-string dichotomy}
% ====================================================================

\subsubsection*{The transgression algebra}

The bosonic string at $c=26$ has a special algebraic structure:
the Koszul dual $\mathrm{Vir}_0$ is uncurved ($\kappa = 0$), so
the bar complex $\barB(\mathrm{Vir}_0)$ is a strict dg coalgebra
($d^2=0$ on the nose).  The cobar construction
$\Omega(\barB(\mathrm{Vir}_0))$ is therefore a strict dga, and
the bar-cobar quasi-isomorphism
$\Omega(\barB(\mathrm{Vir}_0))\xrightarrow{\sim}\mathrm{Vir}_0$
(Theorem~B) is an honest quasi-isomorphism of dgas, not merely of
curved $A_\infty$-algebras.

The \emph{transgression algebra} of $\mathrm{Vir}_c$ at genus~$g$
is a separate construction (\S\ref{subsec:transgression-algebra}):
it is the universal extension $B_\Theta^{(g)} :=
\barB^{(g)}(\mathrm{Vir}_c)\langle\eta\rangle/(\eta^2 -
\kappa\cdot\omega_g)$ that kills the genus-$g$ curvature by
adjoining a formal square root $\eta$ of degree~$1$.  The
extended differential $d_{B_\Theta}^2 = 0$
(Proposition~\ref{prop:transgression-kills-curvature}).

At $c=0$: $\kappa(\mathrm{Vir}_0) = 0$, so $\eta^2 = 0$ and
the transgression algebra degenerates to the exterior extension
$\barB^{(g)} \otimes \bigwedge^*\!\langle\eta\rangle$.  At
$c \neq 0$: $\eta^2 = (c/2)\cdot\omega_g \neq 0$, and the
transgression algebra is a genuine Clifford extension.

The critical-string dichotomy:
\begin{itemize}
\item At $c=0$ (the Koszul dual of $\mathrm{Vir}_{26}$):
  $\kappa = 0$, so $\eta^2 = 0$ and the Clifford relations
  degenerate to exterior algebra relations.  The bar complex is a
  genuine cochain complex; no curved $A_\infty$ machinery is needed.
  The ghost sector is the exterior algebra
  $\bigwedge^*(\Bbbk^{2g})$.
\item At $c\neq 0$: $\eta^2 = (c/2)\omega_g \neq 0$.  The
  transgression algebra is a Clifford extension with
  $\mathrm{Cl}_{2g}(u) \cong \mathrm{Mat}_{2^g}(\Bbbk)$ (Morita
  trivial), and the genus-$g$ structure is controlled by the
  curvature $\kappa = c/2$.
\end{itemize}
The critical dimension $c=26$ is distinguished not by strictness of
the transgression algebra (which is strict at $c=0$, the Koszul
dual side) but by anomaly cancellation: $\kappa_{\mathrm{eff}} =
c/2 + (-13) = 0$ at $c=26$.

\subsubsection*{The coderived passage at $c\neq 26$}

At generic $c \neq 0$, the genus-$g$ bar complex of
$\mathrm{Vir}_c$ is curved ($\kappa(\mathrm{Vir}_c) = c/2 \neq 0$,
so $d_{\barB}^2 = (c/2)\omega_g \neq 0$), which forces
the passage from the derived category to the \emph{coderived}
category.  The coderived category $D^{\mathrm{co}}(\cA^!)$ is
the appropriate triangulated category for curved dg algebras:
it is the Verdier quotient of the homotopy category by the
coacyclic complexes, not the acyclic ones.  The distinction
matters because curved complexes have no cohomology in the
ordinary sense ($d^2\neq 0$), and the correct notion of
``quasi-isomorphism'' must be replaced by ``coacyclic
equivalence.''

The MC element $\Theta_\cA$ at genus $g\ge 1$ naturally lives
in the coderived category: the curvature $\kappa\cdot\omega_g$
is the obstruction to strict nilpotence, and the period
corrections that restore $D^2=0$ on the total space are the
coderived resolution of this obstruction.  The full genus
expansion is a sequence of coderived lifts: each genus-$g$ term
extends the MC solution by one loop order in the coderived
sense.

At $c=26$: the coderived passage is unnecessary for the dual
($\kappa_{\mathrm{dual}}=0$), but the algebra itself still has
curvature $\kappa=13$.  The bosonic string is the theory
where the \emph{dual} is strict but the \emph{algebra} is
still curved.  The critical dimension does not eliminate the
curvature of the boundary theory; it eliminates the curvature
of the bulk.

\subsubsection*{The heuristic BRST/bar comparison}

The proved algebraic input is narrower than a literal
identification of the bosonic-string BRST complex with a bar
construction.  Let $c_{\mathrm{matter}}$ be the matter central
charge and $c_{\mathrm{ghost}}=-26$, so
$c_{\mathrm{total}}=c_{\mathrm{matter}}-26$.  By
Lemma~\ref{lem:brst-nilpotence}, the BRST nilpotence defect is
proportional to $c_{\mathrm{matter}}-26$, and at
$c_{\mathrm{matter}}=26$ the relative BRST complex is an honest
cochain complex.

On the algebraic side, the bar complex
$\barB(\mathrm{Vir}_{c_{\mathrm{total}}})$ has scalar curvature
$\kappa_{c_{\mathrm{total}}}=(c_{\mathrm{matter}}-26)/2$.  Thus
the bar package detects the same anomaly-free locus
$c_{\mathrm{matter}}=26$ at the scalar level.  What is proved
internally is the genus-$0$ total matter+ghost BRST/bar
quasi-isomorphism of Theorem~\ref{thm:brst-bar-genus0}; beyond
that setting, the BRST/bar dictionary is heuristic.  When
$c_{\mathrm{matter}}\neq 26$, the corresponding algebraic
bar/transgression model is curved and naturally belongs to the
curved or coderived formalism, but a physical-state
interpretation still requires additional comparison input.

\subsubsection*{The type-B generalisation}

For type-B strings (worldsheets with boundary), the relevant
algebra is the open-string chiral algebra on the upper
half-plane.  The holographic datum becomes a boundary-bulk-defect
triple $(\cA_{\mathrm{open}},\cA_{\mathrm{closed}},
\cA_{\mathrm{defect}})$, and the transgression algebra acquires
a Swiss-cheese structure.  The critical dimension is unchanged
($c=26$) because the curvature is controlled by the Virasoro
central charge of the closed-string sector.  The open-string
sector contributes the Chan--Paton factors, which modify the
line-operator category but not the curvature.

The Swiss-cheese extension of the modular bar construction
(developed in Volume~II) packages the open-closed-defect system
into a single MC element:
\[
\Theta_{\mathrm{open-closed}}
\;\in\;
\MC\bigl(\mathfrak g^{\mathrm{SC}}_\cA\bigr),
\]
where $\mathfrak g^{\mathrm{SC}}$ is the Swiss-cheese convolution
algebra.  The critical-string dichotomy extends: at $c=26$ the
dual Swiss-cheese algebra is uncurved, and the corresponding
algebraic open-closed transgression model is strict.

\subsubsection*{The superstring analogue}

For the $N=1$ superstring, the boundary chiral algebra is the
$N=1$ super-Virasoro algebra $\mathrm{SVir}_c$ with generators
$T$ (weight~$2$) and $G$ (weight~$3/2$).  The Koszul dual is
$\mathrm{SVir}_{15-c}$, with the super-complementarity
$\kappa(\mathrm{SVir}_c)=(3c-2)/4$
(Remark~\ref{rem:svir-bar-complex}), vanishing at $c=2/3$.
The complementarity sum is
$\kappa(\mathrm{SVir}_c)+\kappa(\mathrm{SVir}_{15-c})=41/4$,
independent of~$c$.  The
critical dimension $c=15$ (ten spacetime dimensions) gives
$\kappa(\mathrm{SVir}_{15})=43/4$.

The self-dual central charge is $c=15/2$: the Koszul involution
acts by $c\mapsto 15-c$, with fixed point $c=15/2$.  The
holographic datum:
\begin{align*}
\mathcal H(\mathrm{SVir}_{15})
&\;=\;
\Bigl(
\mathrm{SVir}_{15},\;\mathrm{SVir}_0,\;\cC_{15},\;
r_{15}(z),\;\Theta_{15},\;\nabla^{\mathrm{hol}}_{15}
\Bigr),\\[4pt]
\kappa_{15}&=\frac{43}{4},
\qquad
\kappa_0=-\frac{1}{2}.
\end{align*}
The pattern parallels the bosonic case: the self-dual point $c=15/2$ is
the midpoint of the Koszul involution, and the shadow obstruction tower
classifies the gravitational complexity.  The fermionic generator~$G$
introduces a super-Mumford correction $\kappa_{\mathrm{ferm}}=(c-2)/4$
that shifts the zero of $\kappa$ from $c=0$ (bosonic) to $c=2/3$ and
gives a nonzero $\kappa_0=-1/2$ at $c=0$.

For the $N=2$ superstring: the boundary algebra is the $N=2$
super-Virasoro $\mathrm{SV}^{(2)}_c$, the critical charge is
$c=9$ (six spacetime dimensions), and the Koszul involution
acts by $c\mapsto 9-c$.  The self-dual point is $c=9/2$.  The
$N=2$ algebra has a spectral-flow automorphism that intertwines
the Ramond and Neveu--Schwarz sectors; this spectral flow is a
symmetry of the holographic datum.

\subsubsection*{The $\mathcal W_\infty$ limit and the higher-spin
critical dimension}

For the $\mathcal W_N$ family, the dual central charge is
$c'_N=K_N-c$ where $K_N$ is the Koszul complementarity
constant ($K_2=26$, $K_3=50$, \dots).  The critical
dimension $c_{\mathrm{crit}}(N)=K_N$ grows with~$N$.  In the
limit $N\to\infty$:
$K_\infty=\lim_{N\to\infty}K_N=\infty$ formally, but the
appropriate statement is that the $\mathcal W_\infty$ algebra has
no finite critical dimension; the dual curvature cannot be
made to vanish at any finite $c$.  The higher-spin critical
dimension is infinite: there is no finite central charge at
which the bulk higher-spin gravity becomes uncurved.

Higher-spin theories in three dimensions are never ``critical''
in the string-theoretic sense: they are always curved.  The
critical dimension $c=26$ is a finite-type phenomenon specific to
the Virasoro algebra; it does not survive the higher-spin limit.

\subsubsection*{Beyond perturbation theory}

The MC element $\Theta_\cA$ is constructed perturbatively (as a
genus expansion), but the bar-intrinsic formula
$\Theta_\cA=D_\cA-d_0$ is exact.  The non-perturbative content
of the holographic datum is encoded in the convergence properties
of the genus expansion.

The HS-sewing theorem guarantees convergence for the standard
landscape.  Beyond the standard landscape (for example, for the
monster vertex algebra $V^\natural$ or for irrational CFTs), the
convergence must be established case by case.  The sewing envelope
$\cA^{\mathrm{sew}}$ provides the analytic framework: a chiral
algebra belongs to the HS class if and only if it embeds into its
sewing envelope with dense image.

The frontier: the three conjectural targets $B_{\mathrm{an}}$
(lattice sewing envelope), $C_{\mathrm{an}}$ (analytic
realisation criterion), $D_{\mathrm{an}}$ (boundary bar duality)
would complete the analytic programme.  The Heisenberg case
(Theorem~\ref{thm:heisenberg-sewing}) is proved; the general
case requires the full strength of the Moriwaki--Adamo--Tanimoto
framework for IndHilb-valued factorisation homology.

\bigskip
\begin{center}
$*\qquad*\qquad*$
\end{center}
\bigskip

\noindent
The gravitational programme extracts no new axioms from the
algebraic engine.  The holographic modular Koszul
datum~\eqref{eq:holo-datum-preface-supp} is a repackaging of
$\Theta_\cA$ and its projections.  The HS-sewing criterion is a
convergence condition on the graph sum.  The $G/L/C/M$
classification is the shadow-depth stratification.  The Koszul
involution $\cA\leftrightarrow\cA^!$ is Verdier duality on
$\Ran(X)$.  The collision $r$-matrix is a residue.  The soft
theorems are Ward identities of the shadow connection.  The
modular bootstrap is the genus spectral sequence.  The critical
dimension is the zero of the dual curvature.

One MC element.  One equation $D^2=0$.  Every gravitational
consequence is a projection.

\bigskip

The eight subsections of this section correspond to eight
chapters of Parts~V--VII of this monograph:

\begin{center}
\small
\renewcommand{\arraystretch}{1.15}
\begin{tabular}{lll}
\textbf{Subsection} & \textbf{Chapter} & \textbf{Key result} \\
\hline
Holographic datum
  & Chapter~\ref{chap:twisted-holography-quantum-gravity}
  & $\mathcal H(T)$ construction \\[2pt]
Perturbative finiteness
  & Theorem~\ref{thm:general-hs-sewing}
  & HS-sewing for standard landscape \\[2pt]
$G/L/C/M$ classes
  & \S\ref{subsec:shadow-postnikov-tower-intro}
  & shadow depth classification \\[2pt]
$S$-duality
  & Theorem C
  & complementarity / Lagrangian \\[2pt]
Gravitational Yangian
  & Chapter~\ref{chap:yangians-foundations}
  & Arnold $\Rightarrow$ CYBE \\[2pt]
Soft theorems
  & \S\ref{sec:modular-shadow-connection}
  & shadow connection Ward identities \\[2pt]
Modular bootstrap
  & \S\ref{sec:genus-spectral-sequence}
  & genus spectral sequence \\[2pt]
Critical-string dichotomy
  & \S\ref{sec:completion-programme-overview}
  & transgression algebra
\end{tabular}
\end{center}

\medskip\noindent
The constructions of Parts~I--III provide the MC element, the
shadow obstruction tower, the Koszul dual, and the genus expansion.  Part~V
interprets these in the gravitational context.  Part~VII pushes
to the frontier.  Volume~II descends to three dimensions.
